
\subsection*{3.3}
\phantomsection
\label{sga3:VIIB:3.3}

If \(G\) is a formal \(k\)-group with affine algebra \(A\), recall
that \(I_A\) denotes the augmentation ideal of \(A\), and that
\[
  \omega_{G/k}
  =
  I_A/I_A^2
  \simeq
  I_A\widehat\otimes_Ak
\]
is a pseudocompact \(k\)-module.

\medskip
\noindent\textbf{Theorem 3.3 \textup{(Cartier)}.}
\phantomsection
\label{sga3:VIIB:3.3:cartier-theorem}
{\itshape Let \(k\) be a pseudocompact ring containing
\(\mathbb Q\), and let \(G\) be a formal \(k\)-group. The following
assertions are equivalent:
\begin{enumerate}
\item[\textup{\SGAInlineRefTarget{sga3:VIIB:3.3:cartier-theorem:item:i}(i)}] There exists an \(\mathcal O_k\)-Lie algebra
\(\mathbf L\), flat over \(\mathcal O_k\), such that \(G\) is
isomorphic to \(\mathscr G(\mathbf L)\). In this case
\(\mathbf L=\operatorname{Lie}(G)\) by \SGARefLink{sga3:VIIB:3.2:prop}{3.2}.

\item[\textup{\SGAInlineRefTarget{sga3:VIIB:3.3:cartier-theorem:item:ii}(ii)}] There exists a projective pseudocompact
\(k\)-module \(\omega\) such that the formal variety underlying
\(G\) is isomorphic to
\[
  \mathbf V_k^{f,0}(\omega)
\]
\textup{(cf.\ \SGARefLink{sga3:VIIB:1.2.5}{1.2.5})}, whose affine algebra is \(k[[\omega]]\).
In this case \(\omega\simeq\omega_{G/k}\).

\item[\textup{\SGAInlineRefTarget{sga3:VIIB:3.3:cartier-theorem:item:iii}(iii)}] The group \(G\) is infinitesimal, and
\(\omega_{G/k}\) is a projective pseudocompact \(k\)-module.

\item[\textup{\SGAInlineRefTarget{sga3:VIIB:3.3:cartier-theorem:item:iv}(iv)}] The group \(G\) is infinitesimal and
topologically flat over \(k\).
\end{enumerate}}

\noindent\textit{Proof of \textup{\SGARefLink{sga3:vii:paragraph:5-1:equation:i}{(i)}}
\(\boldsymbol{\Rightarrow}\) \textup{\SGARefLink{sga3:vii:paragraph:5-1:equation:ii}{(ii)}}.}
Let
\[
  \omega=\Gamma^*(\mathbf L)
\]
be the pseudocompact \(k\)-module dual to \(\mathbf L\)
\textup{(cf.\ \SGARefLink{sga3:VIIB:1.2.3.D}{1.2.3.D})}. For every \(C\in\operatorname{Alf}/k\),
we must construct, functorially in \(C\), an isomorphism
\[
  \mathbf V_k^{f,0}(\omega)(C)
  \xrightarrow{\ \sim\ }
  \mathscr G(\mathbf L)(C).
\]
By \SGARefLink{sga3:VIIB:1.2.5}{1.2.5}, the left-hand side identifies with
\[
  \operatorname{Hom}_c\bigl(\omega,\mathfrak r(C)\bigr),
\]
where \(\mathfrak r(C)\) is the radical of \(C\). This set is
naturally isomorphic to the set
\[
  \operatorname{Hom}_c\bigl(
    \omega\widehat\otimes_kC,\mathfrak r(C)
  \bigr)
\]
of continuous \(C\)-linear maps. Since
\(\omega\widehat\otimes_kC\) is a projective pseudocompact
\(C\)-module, the canonical map
\[
  \bigl(\omega\widehat\otimes_kC\bigr)^\dagger
  \otimes_C\mathfrak r(C)
  \longrightarrow
  \operatorname{Hom}_c\bigl(
    \omega\widehat\otimes_kC,\mathfrak r(C)
  \bigr)
\]
is bijective \textup{(cf.\ \SGARefLink{sga3:VIIB:0.3.6.A}{0.3.6.A})}. By \SGARefLink{sga3:VIIB:1.2.3.E}{1.2.3.E},
\[
  \mathbf L(C)
  \simeq
  \mathbf V_k^f\bigl(\Gamma^*(\mathbf L)\bigr)(C)
  =
  \bigl(\omega\widehat\otimes_kC\bigr)^\dagger.
\]
Thus \(\mathbf V_k^{f,0}(\omega)(C)\) is canonically isomorphic
to
\[
  \mathbf L(C)\otimes_C\mathfrak r(C),
\]
which \SGARefLink{sga3:VIIB:3.2:prop}{Proposition~3.2} identifies canonically with
\(\mathscr G(\mathbf L)(C)\). This proves
\textup{\SGARefLink{sga3:vii:paragraph:5-1:equation:i}{(i)}} \(\Rightarrow\) \textup{\SGARefLink{sga3:VIIB:3.2:prop:item:ii}{(ii)}}.

\medskip
\noindent\textit{Proof of \textup{\SGARefLink{sga3:vii:paragraph:5-1:equation:ii}{(ii)}}
\(\boldsymbol{\Rightarrow}\) \textup{\SGARefLink{sga3:VIIB:3.3:cartier-theorem:item:iii}{(iii)}}.}
Let \(\omega\) be a projective object of \(\operatorname{PC}(k)\),
and let
\[
  h:k[[\omega]]\xrightarrow{\ \sim\ }A
\]
be an isomorphism onto the affine algebra of \(G\). The composite
\[
  \varepsilon_A\circ h:k[[\omega]]\longrightarrow k
\]
is determined by its restriction
\(\lambda:\omega\to\mathfrak r(k)\). The map
\[
  x\longmapsto x-\lambda(x)
\]
from \(\omega\) to the radical of \(k[[\omega]]\) extends, by the
universal property of \(k[[\omega]]\), to an endomorphism
\(\ell_\lambda\) of \(k[[\omega]]\). Since
\[
  \ell_\lambda\circ\ell_{-\lambda}
  =
  \ell_{-\lambda}\circ\ell_\lambda
  =
  \operatorname{id},
\]
the map \(\ell_\lambda\) is an automorphism. Replacing \(h\) by
\(h\circ\ell_\lambda\), we may assume that
\(\varepsilon_A\circ h\) vanishes on the closed ideal \(I\) of
\(k[[\omega]]\) generated by \(\omega\). Then \(h\) induces an
isomorphism
\[
  I/I^2\xrightarrow{\ \sim\ }I_A/I_A^2.
\]
Since \(I/I^2\simeq\omega\), we obtain
\[
  \omega_{G/k}\simeq\omega,
\]
so \(\omega_{G/k}\) is projective. The formal variety
\(\mathbf V_k^{f,0}(\omega)\), and hence \(G\), is infinitesimal.

\medskip
\noindent\textit{Proof of \textup{\SGARefLink{sga3:VIIB:3.3:cartier-theorem:item:iii}{(iii)}}
\(\boldsymbol{\Rightarrow}\) \textup{\SGARefLink{sga3:vii:paragraph:5-1:equation:i}{(i)}}.}
Assume that \(G\) is infinitesimal and \(\omega_{G/k}\) is
projective. Let \(\mathbf L\) be the \(\mathcal O_k\)-Lie algebra
of \(G\). By \SGARefLink{sga3:VIIB:2.6.1}{2.6.1}, its underlying \(\mathcal O_k\)-module is
\[
  \mathbf V_k^f(\omega_{G/k}).
\]
Consequently, \(\mathbf L\) is flat over \(\mathcal O_k\), and
\SGARefLink{sga3:VIIB:1.2.3.E}{Proposition~1.2.3.E} gives
\[
  \Gamma^*(\mathbf L)\simeq\omega_{G/k}.
\]
The proof of \textup{\SGARefLink{sga3:vii:paragraph:5-1:equation:i}{(i)}} \(\Rightarrow\) \textup{\SGARefLink{sga3:vii:paragraph:5-1:equation:ii}{(ii)}} therefore
identifies the affine algebra of \(\mathscr G(\mathbf L)\) with
\(k[[\omega_{G/k}]]\). By \SGARefLink{sga3:VIIB:2.6.3:prop}{2.6.3}, the identity morphism of
\(\mathbf L\) corresponds canonically to a morphism
\[
  \mathscr G(\mathbf L)\longrightarrow G,
\]
and hence to a continuous \(k\)-algebra morphism
\[
  \varphi:A\longrightarrow k[[\omega_{G/k}]].
\]
Let \(I\) be the closed ideal generated by \(\omega_{G/k}\).
Filter \(k[[\omega_{G/k}]]\) by the closures of the powers \(I^n\),
and \(A\) by the closures of the powers \(I_A^n\). By definition,
\(\varphi\) induces the identity on
\[
  \omega_{G/k}=I_A/I_A^2=I/I^2.
\]

Because \(\omega_{G/k}\) is projective, the canonical projection
\[
  I_A\longrightarrow\omega_{G/k}
\]
has a section \(\tau\). By the universal property of
\(k[[\omega_{G/k}]]\), this section defines a continuous algebra
morphism
\[
  \psi:k[[\omega_{G/k}]]\longrightarrow A.
\]
The composite \(\varphi\circ\psi\) induces the identity on
\(\omega_{G/k}\), and hence on the associated graded algebra of
\(k[[\omega_{G/k}]]\). It follows from
\textup{[CA], \S V.7, Lemma~1} that \(\varphi\circ\psi\) is an
isomorphism.\footnote{\label{sga3:VIIB:nde:126}\textit{Editors' note.} Compare \SGARefLink{sga3:VIIB:5.1.5}{5.1.5}
below, and also \textup{[BAC], III, \S \SGARefLink{sga3:VIIB:2.8}{2.8}, Theorem~1 and its
corollaries}. The editors supplied the details of the original
argument that follow.}

Moreover, \(\psi\) induces an isomorphism
\[
  I/I^2\xrightarrow{\ \sim\ }I_A/I_A^2,
\]
and hence a surjection on associated graded algebras. Since \(G\)
is radicial, \(I_A\) is contained in the radical of \(A\), so
\[
  \bigcap_n I_A^n=0.
\]
The cited lemma therefore shows that \(\psi\) is surjective. Since
\(\varphi\circ\psi\) is an isomorphism and \(\psi\) is surjective,
both \(\psi\) and \(\varphi\) are isomorphisms. This proves
\textup{\SGARefLink{sga3:VIIB:3.3:cartier-theorem:item:iii}{(iii)}} \(\Rightarrow\) \textup{\SGARefLink{sga3:VIIB:3.3:cartier-theorem:item:i}{(i)}}.

\medskip
Finally, \textup{\SGARefLink{sga3:vii:paragraph:5-1:equation:i}{(i)}} or \textup{\SGARefLink{sga3:vii:paragraph:5-1:equation:ii}{(ii)}} plainly implies
\textup{\SGARefLink{sga3:VIIB:3.3:cartier-theorem:item:iv}{(iv)}}, so it remains to prove
\textup{\SGARefLink{sga3:VIIB:3.3:cartier-theorem:item:iv}{(iv)}} \(\Rightarrow\) \textup{\SGARefLink{sga3:VIIB:3.3:cartier-theorem:item:ii}{(ii)}}. We may assume that
\(k\) is local, with residue field \(k_0\). Put
\[
  G_0=G\widehat\otimes_kk_0,
  \qquad
  \omega=\omega_{G/k},
  \qquad
  \omega_0=\omega\widehat\otimes_kk_0.
\]
\footnote{\label{sga3:VIIB:nde:127}The editors supplied the
details of the original argument that follow.}
The pseudocompact \(k_0\)-module \(\omega_0\) has a pseudobasis
\((y_\lambda)_{\lambda\in\Lambda}\). Let
\[
  \omega'=\prod_{\lambda\in\Lambda}k_\lambda
\]
be the topologically free \(k\)-module with that index set, and
lift each \(y_\lambda\) to an element \(x_\lambda\in\omega\).
This gives a continuous \(k\)-linear map
\[
  f:\omega'\longrightarrow\omega
\]
such that
\[
  f_0=f\widehat\otimes_kk_0
\]
is invertible.
Since \(\omega'\) is projective, \(f\) lifts to a continuous
\(k\)-linear map
\[
  g:\omega'\longrightarrow I_A
\]
with \(\pi\circ g=f\), where
\(\pi:I_A\to\omega\) is the projection. The universal property
of \(k[[\omega']]\) gives a morphism of topological algebras
\[
  \phi:k[[\omega']]\longrightarrow A.
\]

Now
\[
  \omega'\widehat\otimes_kk_0
  \simeq
  \omega\widehat\otimes_kk_0
  =
  \omega_{G_0/k_0},
\]
and consequently
\[
  k[[\omega']]\widehat\otimes_kk_0
  \simeq
  k_0[[\omega_{G_0/k_0}]].
\]
The pair \(k_0,G_0\) satisfies \textup{\SGARefLink{sga3:VIIB:3.3:cartier-theorem:item:iii}{(iii)}}, and the proof of
\textup{\SGARefLink{sga3:VIIB:3.3:cartier-theorem:item:iii}{(iii)}} \(\Rightarrow\) \textup{\SGARefLink{sga3:vii:paragraph:5-1:equation:i}{(i)}} shows that
\[
  \phi_0=\phi\widehat\otimes_kk_0
\]
is invertible. Both \(k[[\omega']]\) and \(A\) are projective
pseudocompact \(k\)-modules, so \SGARefLink{sga3:VIIB:0.3.4}{0.3.4} implies that \(\phi\) itself
is invertible. In particular, if \(I\) is the augmentation ideal
of \(k[[\omega']]\), then \(\phi\) induces an isomorphism
\[
  \omega'=I/I^2
  \xrightarrow{\ \sim\ }
  I_A/I_A^2=\omega.
\]
This completes the proof.
